\scanpage{80}

\deleted{\deleted{\unsure{Prouvons|Réunissons}\,;} Il suffit de prouver que \deleted{\unsure{tout} \illegible{2} \unsure{possible} $\sum\unsure{(z_1+\cdots+z_n)}$ est \unsure{majoré} par $\|u\|_1$ \unsure{distinct}. Soit \unsure{à cet effet}} $x_1,\ldots,x_n$ des vecteurs propres lin. \deleted{\unsure{à des valeurs propres}} indépendants associés aux valeurs propres $z_i$, \deleted{\unsure{soit}} on a $\sum|z_i|\leq\|u\|_1$. Soit en effet}

% margin note written sideways
\margin{\deleted{Nous prouvons au N° suivant que $\sum|\lambda_i|\leq\|u\|_1$. Ici, nous nous bornons à établir, plus élémentairement et \deleted{\unsure{par}} une proposition \unsure{pourtant plus faible},}}

\deleted{\deleted{\unsure{soit}} $p$ l'opérateur de projection orthogonale de $E$ sur le sous espace engendré par les $x_i$, on a}

\[
\deleted{|z_1+\cdots+z_n|=|\mathrm{Tr}\,pu|\leq\|pu\|_1\leq\|p\|\,\|u\|_1=\|u\|_1.}
\]

\deleted{\deleted{Comme} Les \deleted{\unsure{zéros}} inverses des zéros de $\det(1-zu)$ sont les $z_i$, et forment donc une suite sommable, que d'autre part les coefficients de $\det(1-zu)$ sont majorés, terme à terme, en valeur, par ceux de $\det(1+zh)$ où $h=\sqrt{u^*u}$ (prop. 1 corollaire 4). Nous supposons connu qu'une fonction entière dont les inverses des zéros forment une suite sommable, et \deleted{\unsure{que}} dont les coefficients sont majorés en module par ceux d'une série entière de genre 0, est elle \unsure{=|même} de genre 0. \deleted{\unsure{Ceci étant}, Prouvons} Il reste \deleted{\unsure{donc}}\inserted{donc} seulement à montrer que $\det(1+zh)$ est de genre 0.}

\deleted{\deleted{\unsure{Commençons} \inserted{donc} à vérifier les formules (23) pour $u=h$ \unsure{hermitien} $\geq0$, \unsure{les valeurs propres} \illegible{1} positives. Mais \deleted{\unsure{ce sont alors}} l'identité (22) signifie que $(\unsure{3|23|22})$ $\alpha_n(u)=\sum_{i_1<\cdots<i_n}\lambda_{i_1}\ldots\lambda_{i_n}$, ce qui se vérifie en effet immédiatement sur l'expression}}

\deleted{On aura $u=\sum\lambda_i\bar e_i\otimes e_i$ (où $(e_i)$ est un système orthonormal), d'où \deleted{\unsure{En effet, on aura}}\deleted{\unsure{déjà|donc}}}

\[
\deleted{\alpha_n(u)=\sum_{i_1,\ldots,i_n}\deleted{\unsure{\alpha_n(z_1e_1,\ldots)|\alpha_n(z_i e,\ldots)}}z_{i_1}\cdots z_{i_n}\,\alpha_n(\bar e_{i_1}\otimes e_{i_1},\ldots,\bar e_{i_n}\otimes e_{i_n}).}
\]

\deleted{\deleted{\unsure{et}} \deleted{\unsure{et}} $\alpha_n(\bar e_{i_1}\otimes e_{i_1},\ldots,\bar e_{i_n}\otimes e_{i_n})$ \deleted{\unsure{est}} est nul, sauf quand $i_1,\ldots,i_n$ sont distincts, auquel cas c'est égal à $\frac1{n!}$\,; rassemblant dans la somme les termes pour lesquels les $i_1,\ldots,i_n$ ne diffèrent que par l'ordre, on trouve bien (23).}

% archive sheet number: 81
